
% =========================================================
\section{Smooth flip bifurcation and the emerging period-two branch}
\label{sec:flip-bifurcation}
% =========================================================

In this section we study the smooth gradient map
\[
G_\eta(x)=x-\eta\nabla V(x)
\]
near the critical step size
\[
\eta_f=\frac{2}{\lambda_*}.
\]

The local mechanism considered here is the classical simple flip
bifurcation of a smooth discrete map
\cite{Kuznetsov2004,GuckenheimerHolmes1983}.
For gradient-descent dynamics, Mulayoff and Stich
\cite{MulayoffStich2026} recently derived a multivariate nonlinear
stability criterion involving the second-, third-, and fourth-order
derivatives of the potential, while Gan \cite{Gan2026} formulated
edge-of-stability dynamics in a systematic bifurcation-theoretic
framework.

Accordingly, the existence of a smooth period-two branch is not the
new contribution of the present paper.  We reproduce the local
reduction because its explicit coefficient enters the
switching-separation law derived later; we therefore give a
self-contained center-manifold and normal-form derivation in the
notation of the present paper.

By the local identity theorem established in Section~3, the
resulting local bifurcation structure is inherited exactly by the
clipped map $F_{\eta,c}$ as long as the bifurcating orbit remains
inside the smooth region $\Omega_c^-$.

Throughout this section we write
\begin{equation}
\mu
=
\eta-\eta_f.
\label{eq:flip-mu}
\end{equation}
Thus the bifurcation point corresponds to $\mu=0$.

Under \textbf{(H2)--(H3)}, the multiplier associated with $v_*$ is
\begin{equation}
\rho_*(\mu)
=
1-(\eta_f+\mu)\lambda_*
=
-1-\lambda_*\mu,
\label{eq:critical-multiplier-mu}
\end{equation}
whereas all remaining multipliers lie strictly inside the unit disk
at $\mu=0$.

The nonlinear character of the bifurcation is governed by the
coefficient
\[
\ell_*
=
\frac{\Gamma_*}{3\lambda_*},
\]
defined in \eqref{eq:ell-star}.  Assumption \textbf{(H4)} gives
\[
\ell_*>0,
\]
which corresponds to the supercritical case.

% ---------------------------------------------------------
\subsection{Parameter-dependent center reduction}
\label{subsec:parameter-center-reduction}
% ---------------------------------------------------------

Translate the critical point to the origin, $y=x-x_*$, and define
\begin{equation}
\Phi_\mu(y)
:=
G_{\eta_f+\mu}(x_*+y)-x_*.
\label{eq:Phi-mu}
\end{equation}
Then $\Phi_\mu(0)=0$ for every sufficiently small $\mu$, and at
$\mu=0$,
\begin{equation}
D\Phi_0(0)
=
I-\eta_fH_*,
\label{eq:Phi-linearization}
\end{equation}
whose unique critical multiplier is $-1$ with eigenvector $v_*$; the
complementary spectrum lies strictly inside the unit disk.

To incorporate the bifurcation parameter, consider the extended map
\begin{equation}
(y,\mu)
\longmapsto
\bigl(\Phi_\mu(y),\mu\bigr).
\label{eq:extended-map}
\end{equation}
Its two-dimensional local center manifold is associated with the
critical state direction $v_*$ and the neutral parameter direction,
and after restricting to a sufficiently small neighborhood can be
represented as
\begin{equation}
y
=
zv_*+h(z,\mu),
\label{eq:center-parametrization}
\end{equation}
with
\begin{equation}
h(z,\mu)\in v_*^\perp,
\qquad
h(0,\mu)=0,
\qquad
D_zh(0,0)=0,
\qquad
h(z,\mu)
=
O\bigl(z^2+|\mu||z|\bigr).
\label{eq:center-h-properties}
\end{equation}
The dynamics on the center manifold are described by the scalar map
\begin{equation}
z_{n+1}
=
f_\mu(z_n).
\label{eq:reduced-map}
\end{equation}

% ---------------------------------------------------------
\subsection{Flip normal form}
\label{subsec:flip-normal-form}
% ---------------------------------------------------------

We next record the local normal form required for the asymptotic
analysis.

\begin{proposition}[Reduced flip normal form]
\label{prop:flip-normal-form}
Assume \textbf{(H1)--(H4)}.
There exist local state and parameter coordinates, preserving
$\mu=0$ and $z=0$, in which the reduced map
\eqref{eq:reduced-map} takes the form
\begin{equation}
f_\mu(z)
=
-\bigl(1+\lambda_*\mu\bigr)z
+
\ell_*z^3
+
R(z,\mu),
\label{eq:flip-normal-form}
\end{equation}
where
\begin{equation}
R(z,\mu)
=
O\left(
|z|^4
+
|\mu||z|^2
+
\mu^2|z|
\right)
\label{eq:flip-remainder}
\end{equation}
as $(z,\mu)\to(0,0)$.

The cubic coefficient is
\begin{equation}
\ell_*
=
\frac{\eta_f}{6}
\left(
3T_*[v_*,v_*,H_*^{-1}g_*]
-
Q_*[v_*,v_*,v_*,v_*]
\right),
\label{eq:flip-coefficient-section3}
\end{equation}
or equivalently
\begin{equation}
\ell_*
=
\frac{\Gamma_*}{3\lambda_*}.
\label{eq:flip-coefficient-Gamma}
\end{equation}
\end{proposition}

\begin{proof}
The critical multiplier $-1$ is simple and crosses the unit circle
transversally (Section~\ref{subsec:spectral-consequences}), so the
standard parameter-dependent center-manifold and flip-normal-form
reductions apply; see \cite{Carr1981,HirschPughShub1977,Kuznetsov2004}.
The linear parameter coefficient is $\rho_*(\mu)=-1-\lambda_*\mu$
by \eqref{eq:critical-multiplier-mu}, and the quadratic term is
eliminated by the standard near-identity transformation for a simple
flip; the remaining cubic coefficient is the first flip coefficient.

For the gradient map
\[
G_\eta(x)=x-\eta\nabla V(x),
\]
one has
\begin{align}
D^2G_{\eta_f}(x_*)[u,w]
&=
-\eta_f
T_*[u,w,\cdot]^\sharp,
\label{eq:D2G}
\\
D^3G_{\eta_f}(x_*)[u,w,q]
&=
-\eta_f
Q_*[u,w,q,\cdot]^\sharp.
\label{eq:D3G}
\end{align}
The standard flip-coefficient formula therefore gives
\[
\ell_*
=
\frac{\eta_f}{2}
T_*[v_*,v_*,H_*^{-1}g_*]
-
\frac{\eta_f}{6}
Q_*[v_*,v_*,v_*,v_*].
\]
Using
\[
\eta_f=\frac{2}{\lambda_*}
\]
and the definition of $\Gamma_*$ in
\eqref{eq:Gamma-star} yields
\[
\ell_*
=
\frac{\Gamma_*}{3\lambda_*}.
\]

The remainder estimate
\eqref{eq:flip-remainder} follows from the regularity in
\textbf{(H1)} and the parameter-dependent near-identity normal-form
transformation.
\end{proof}

\begin{remark}
\label{rem:flip-sign-convention}
The sign convention in \eqref{eq:flip-normal-form} is fixed
throughout the paper:
\[
\ell_*>0
\]
corresponds to a supercritical flip bifurcation, with a stable
period-two branch emerging for $\mu>0$.
This convention is consistent with the multivariate
gradient-descent flip criterion encoded by \textbf{(H4)}.
\end{remark}

% ---------------------------------------------------------
\subsection{The second iterate}
\label{subsec:second-iterate}
% ---------------------------------------------------------

The bifurcating period-two orbit is most conveniently obtained from
the fixed points of the second iterate.

\begin{lemma}[Expansion of the second iterate]
\label{lem:second-iterate-expansion}
Under the hypotheses of
Proposition~\ref{prop:flip-normal-form},
\begin{equation}
f_\mu^2(z)-z
=
2\lambda_*\mu z
-
2\ell_*z^3
+
\mathcal R_2(z,\mu),
\label{eq:second-iterate-expansion}
\end{equation}
where
\begin{equation}
\mathcal R_2(z,\mu)
=
O\left(
|z|^4
+
|\mu||z|^2
+
\mu^2|z|
\right).
\label{eq:second-iterate-remainder}
\end{equation}
\end{lemma}

\begin{proof}
Write
\[
a_\mu:=1+\lambda_*\mu.
\]
Then
\[
f_\mu(z)
=
-a_\mu z+\ell_*z^3+R(z,\mu).
\]
Ignoring terms of order contained in
\eqref{eq:second-iterate-remainder}, one obtains
\begin{align}
f_\mu^2(z)
&=
-a_\mu
\left(
-a_\mu z+\ell_*z^3
\right)
+
\ell_*
\left(
-a_\mu z
\right)^3
+
\mathcal R_2(z,\mu)
\nonumber\\
&=
a_\mu^2z
-
\ell_*a_\mu(1+a_\mu^2)z^3
+
\mathcal R_2(z,\mu).
\label{eq:second-compose}
\end{align}
Since
\[
a_\mu^2
=
1+2\lambda_*\mu+O(\mu^2)
\]
and
\[
a_\mu(1+a_\mu^2)
=
2+O(\mu),
\]
we obtain
\[
f_\mu^2(z)-z
=
2\lambda_*\mu z
-
2\ell_*z^3
+
\mathcal R_2(z,\mu).
\]
The stated remainder estimate follows by substituting
\eqref{eq:flip-remainder} into the composition.
\end{proof}

Equation \eqref{eq:second-iterate-expansion} exhibits explicitly the
balance responsible for the square-root amplitude law:
\[
\lambda_*\mu
\sim
\ell_*z^2.
\]

% ---------------------------------------------------------
\subsection{Existence and amplitude of the period-two branch}
\label{subsec:period-two-amplitude}
% ---------------------------------------------------------

We now obtain the local period-two branch together with its leading
amplitude.

\begin{theorem}[Supercritical period-two branch]
\label{thm:period-two-branch}
Assume \textbf{(H1)--(H4)}.
There exists $\mu_0>0$ such that for every
\[
0<\mu<\mu_0,
\]
the smooth gradient map $G_{\eta_f+\mu}$ possesses a unique local
period-two cycle
\begin{equation}
\mathcal O_2(\mu)
=
\left\{
x_-(\mu),x_+(\mu)
\right\}
\label{eq:period-two-cycle}
\end{equation}
in a sufficiently small neighborhood of $x_*$.

The two points satisfy
\begin{equation}
x_\pm(\mu)
=
x_*
\pm
A_*\mu^{1/2}v_*
+
O(\mu),
\qquad
\mu\to0^+,
\label{eq:period-two-x-expansion}
\end{equation}
where
\begin{equation}
A_*
=
\sqrt{\frac{\lambda_*}{\ell_*}}
\label{eq:A-star-ell}
\end{equation}
and therefore
\begin{equation}
A_*^2
=
\frac{3\lambda_*^2}{\Gamma_*}.
\label{eq:A-star-Gamma}
\end{equation}

The labeling may be chosen so that
\begin{equation}
G_{\eta_f+\mu}(x_+(\mu))
=
x_-(\mu),
\qquad
G_{\eta_f+\mu}(x_-(\mu))
=
x_+(\mu).
\label{eq:period-two-exchange}
\end{equation}
\end{theorem}

\begin{proof}
We first work in the reduced coordinate $z$.

Set
\[
\mu=\varepsilon^2,
\qquad
\varepsilon>0,
\]
and seek nonzero solutions of
\[
f_{\varepsilon^2}^2(z)-z=0
\]
in the form
\[
z=\varepsilon w.
\]
By Lemma~\ref{lem:second-iterate-expansion},
\begin{align}
f_{\varepsilon^2}^2(\varepsilon w)-\varepsilon w
&=
2\lambda_*\varepsilon^3w
-
2\ell_*\varepsilon^3w^3
+
O(\varepsilon^4)
\nonumber\\
&=
2\varepsilon^3
\left[
w(\lambda_*-\ell_*w^2)
+
O(\varepsilon)
\right].
\label{eq:scaled-second-iterate}
\end{align}

Define the rescaled bifurcation equation
\begin{equation}
\mathscr F(w,\varepsilon)
:=
\frac{
f_{\varepsilon^2}^2(\varepsilon w)-\varepsilon w
}{
2\varepsilon^3
},
\qquad
\varepsilon>0.
\label{eq:rescaled-F}
\end{equation}
Under \textbf{(H1)} the rescaled map extends as a $C^3$ function to
$\varepsilon=0$: since $V\in C^7$, the gradient map $G_\eta$ is
$C^6$, and by the parameter-dependent flip reduction so is the
reduced map $f_\mu$; the numerator in \eqref{eq:rescaled-F} is
therefore $C^6$ in $(\varepsilon,w)$ and vanishes to third order in
$\varepsilon$, so dividing by $2\varepsilon^3$ yields a $C^3$
function.  We set
\begin{equation}
\mathscr F(w,0)
=
w(\lambda_*-\ell_*w^2).
\label{eq:rescaled-limit}
\end{equation}

Since $\ell_*>0$, the nonzero zeros of
$\mathscr F(\cdot,0)$ are
\begin{equation}
w_\pm^0
=
\pm
\sqrt{\frac{\lambda_*}{\ell_*}}.
\label{eq:w0}
\end{equation}
Furthermore,
\begin{align}
\partial_w\mathscr F(w_\pm^0,0)
&=
\lambda_*
-
3\ell_*(w_\pm^0)^2
\nonumber\\
&=
-2\lambda_*
\neq0.
\label{eq:IFT-nondegenerate}
\end{align}
The implicit function theorem therefore yields two branches
\[
w_\pm(\varepsilon)
=
w_\pm^0+O(\varepsilon).
\]
Consequently,
\begin{align}
z_\pm(\mu)
&=
\varepsilon w_\pm(\varepsilon)
\nonumber\\
&=
\pm
\sqrt{\frac{\lambda_*}{\ell_*}}
\sqrt{\mu}
+
O(\mu).
\label{eq:z-period-two}
\end{align}

The fixed point $z=0$ is locally the only fixed point of
$f_\mu$ near the origin.  Indeed,
\[
\partial_z(f_0(z)-z)\big|_{z=0}
=
-2
\neq0,
\]
so the implicit function theorem excludes nearby nonzero
period-one fixed points.  Hence the two nonzero solutions
$z_\pm(\mu)$ of
\[
f_\mu^2(z)=z
\]
have prime period two.

Returning to the original coordinates through
\eqref{eq:center-parametrization},
\[
x_\pm(\mu)
=
x_*
+
z_\pm(\mu)v_*
+
h(z_\pm(\mu),\mu).
\]
Since
\[
h(z,\mu)
=
O(z^2+|\mu||z|)
\]
and
\[
z_\pm(\mu)=O(\sqrt{\mu}),
\]
we have
\[
h(z_\pm(\mu),\mu)=O(\mu).
\]
Therefore
\[
x_\pm(\mu)
=
x_*
\pm
\sqrt{\frac{\lambda_*}{\ell_*}}
\sqrt{\mu}\,v_*
+
O(\mu),
\]
which proves \eqref{eq:period-two-x-expansion}.

Finally, the two nonzero points form a single two-cycle, so after
possibly interchanging the labels we obtain
\eqref{eq:period-two-exchange}.

The uniqueness of the sufficiently small period-two cycle follows
from the local center-manifold reduction and the uniqueness of the
two nonzero solution branches of the reduced second-iterate
equation.
\end{proof}

% ---------------------------------------------------------
\subsection{Stability of the emerging cycle}
\label{subsec:period-two-stability}
% ---------------------------------------------------------

We next verify the local stability of the branch.

\begin{theorem}[Stability of the supercritical period-two cycle]
\label{thm:period-two-stability}
Under \textbf{(H1)--(H4)}, the period-two cycle
$\mathcal O_2(\mu)$ given by
Theorem~\ref{thm:period-two-branch} is locally asymptotically stable
for all sufficiently small $\mu>0$.

Along the critical center direction, the multiplier of the
two-step map satisfies
\begin{equation}
\rho_{\mathrm{cyc}}^{c}(\mu)
=
1-4\lambda_*\mu
+
O(\mu^{3/2}),
\qquad
\mu\to0^+.
\label{eq:critical-cycle-multiplier}
\end{equation}
\end{theorem}

\begin{proof}
Differentiate the expansion
\eqref{eq:second-iterate-expansion}:
\begin{equation}
(f_\mu^2)'(z)
=
1
+
2\lambda_*\mu
-
6\ell_*z^2
+
O\left(
|z|^3
+
|\mu||z|
+
\mu^2
\right).
\label{eq:second-iterate-derivative}
\end{equation}
At either point $z=z_\pm(\mu)$,
Theorem~\ref{thm:period-two-branch} gives
\[
z_\pm(\mu)^2
=
\frac{\lambda_*}{\ell_*}\mu
+
O(\mu^{3/2}).
\]
Hence
\begin{align}
(f_\mu^2)'(z_\pm(\mu))
&=
1
+
2\lambda_*\mu
-
6\lambda_*\mu
+
O(\mu^{3/2})
\nonumber\\
&=
1
-
4\lambda_*\mu
+
O(\mu^{3/2}).
\end{align}
Thus the center multiplier lies strictly inside the unit disk for
all sufficiently small $\mu>0$.

At $\mu=0$, every transverse multiplier of the two-step map equals
\[
\rho_i(\eta_f)^2,
\qquad
i=1,\ldots,d-1,
\]
and by \eqref{eq:noncritical-spectrum},
\[
|\rho_i(\eta_f)|<1.
\]
Therefore
\[
\rho_i(\eta_f)^2<1.
\]
By continuity, all transverse multipliers remain strictly inside the
unit disk along the sufficiently small period-two branch.

The reduction principle then implies local asymptotic stability of
the full period-two cycle.
\end{proof}

\begin{remark}
\label{rem:cycle-multiplier-meaning}
Equation \eqref{eq:critical-cycle-multiplier} shows that the
newly created cycle is only weakly attracting immediately after the
flip:
\[
\rho_{\mathrm{cyc}}^{c}(\mu)\to1^-
\qquad
\text{as }\mu\to0^+.
\]
This weak attraction is relevant later when the branch approaches
the switching hypersurface, because the switching interaction occurs
on a parameter scale that also tends to zero with $c$.
\end{remark}

% ---------------------------------------------------------
\subsection{Transfer to the clipped map before switching}
\label{subsec:transfer-clipped}
% ---------------------------------------------------------

We finally transfer the preceding smooth bifurcation result to the
clipped map.

\begin{corollary}[Pre-contact period-two cycle for the clipped map]
\label{cor:precontact-clipped-cycle}
Assume \textbf{(H1)--(H5)} and fix $c>0$.
Then there exists
\[
\mu_c>0
\]
such that for every
\[
0<\mu<\mu_c,
\]
the clipped map
\[
F_{\eta_f+\mu,c}
\]
possesses the same locally asymptotically stable period-two cycle
$\mathcal O_2(\mu)$ as the smooth gradient map
$G_{\eta_f+\mu}$.

Moreover,
\begin{equation}
x_\pm(\mu)
\in
\Omega_c^-,
\qquad
0<\mu<\mu_c.
\label{eq:cycle-inside-smooth-region}
\end{equation}
\end{corollary}

\begin{proof}
For the fixed value $c>0$, let $r_c>0$ be given by
Theorem~\ref{thm:local-identity}. By
Theorem~\ref{thm:period-two-branch},
\[
x_\pm(\mu)\to x_*
\qquad
\text{as }\mu\to0^+.
\]
Hence there exists $\mu_c>0$ such that
\[
x_\pm(\mu)\in B_{r_c}(x_*)
\subset\Omega_c^-
\]
for every $0<\mu<\mu_c$.

On this ball,
\[
F_{\eta_f+\mu,c}
=
G_{\eta_f+\mu}
\]
by Theorem~\ref{thm:local-identity}. Therefore the two maps have
exactly the same period-two cycle and the same local derivatives
along that cycle.  Its local asymptotic stability follows from
Theorem~\ref{thm:period-two-stability}.
\end{proof}

\begin{remark}[The first possible effect of clipping]
\label{rem:first-clipping-effect}
Corollary~\ref{cor:precontact-clipped-cycle} shows that the
period-two branch is born entirely inside the smooth region.
Therefore clipping does not influence either its creation or its
initial stability.

The first possible effect of clipping occurs only when one of the
points $x_\pm(\mu)$ reaches the endogenous switching hypersurface
\[
\Sigma_c
=
\left\{
x:
\|\nabla V(x)\|_\diamond=c
\right\}.
\]
Determining the parameter value of this first contact is the central
problem of Sections~\ref{sec:switching-geometry} and
\ref{sec:separation-law}.
\end{remark}

